\subsection{Discussion}\label{sec:exp-discussion}

\paragraph{$\sigma_\ell$ is a diagnostic, not a diversity signal.} $\sigma_\ell$ increased monotonically with $m$ in our synthetic experiments because our modes have very different formats (code vs.\ haiku vs.\ recipe $\to$ different per-byte cross-entropies).
But $\sigma_\ell$ measures coherence \emph{heterogeneity}, not diversity: modes with similar fluency would have low $\sigma_\ell$ despite high diversity, and a single mode with high quality variance could have high $\sigma_\ell$ with zero diversity.
Its role is diagnostic (detecting mixed coherence), not as a standalone diversity score.

\paragraph{Cross-mode information scales with model quality.} We test how cross-mode information transfer varies with $\theta$'s scale by running the pairwise matrix experiment on \scalingNumLlamaModels{} dense Llama models (1B, 3B, 8B, 70B) using the same \crossmodeNumModes{} modes and 5 samples per mode.
Figure~\ref{fig:scaling-crossmode} shows the results.
The off-diagonal mean increases monotonically with model size: $\scalingOffDiagLlamaOneB$ (1B), $\scalingOffDiagLlamaThreeB$ (3B), $\scalingOffDiagLlamaEightB$ (8B), $\scalingOffDiagLlamaSeventyB$ (70B) bits, transitioning from negative (cross-mode context hurts) to positive (cross-mode context helps).
The fraction of positive off-diagonal entries follows the same monotone trend: \scalingFracPosLlamaOneB\%, \scalingFracPosLlamaThreeB\%, \scalingFracPosLlamaEightB\%, \scalingFracPosLlamaSeventyB\%.
This hypothesis was pre-registered before running any Llama models, and the probability of accidental monotonicity across \scalingNumLlamaModels{} models under the null is $1/\scalingNumLlamaModels! \approx \scalingFactorialProbPct\%$.

\begin{figure*}[t]
\centering
\includegraphics[width=\textwidth]{scaling_comparison/paper_scaling_figure.pdf}
\caption{Cross-mode information transfer scales with model size. \textbf{Left}: Mean off-diagonal surprise reduction (with bootstrap 95\% CI) transitions from negative to positive as models grow, indicating that larger models extract more information from cross-mode context. \textbf{Right}: The fraction of mode pairs showing positive cross-mode transfer increases monotonically across the \scalingNumLlamaModels{} Llama models (blue; $p = \scalingFactorialProbPct\%$ under monotone null).
GPT-2 and Qwen2.5-3B (gray) are included for context but differ in architecture.}
\label{fig:scaling-crossmode}
\end{figure*}

\paragraph{The framework admits other metrics.} The $a_k$ curve is the primary object of our framework, and $D_{Ca_n} = C \times a_n$ is one natural summary of it among many.
Reporting the curve directly preserves information that any single scalar (including ours) collapses.
Alternative scalars could be derived from the same curve to capture different desiderata: weighted sums of $a_k - a_\infty$ (emphasizing early or late positions), the slope or curvature at a specific $k$ (measuring how quickly $\theta$ learns), curve-shape descriptors, or coherence-weighted variants that combine $C$ with quantities other than $a_n$.
Different applications may warrant different choices; for instance, a policy-comparison setting might prefer a metric sensitive to \emph{changes} in the curve shape, while a ranking setting needs a single well-behaved scalar.
We have not explored this space systematically; we focus on $D_{Ca_n}$ because it empirically works well across regimes and has a clean information-theoretic interpretation, but we expect the framework to support a family of related metrics with complementary properties.

